\documentclass[11pt,a4paper]{article}

% --- Mathematics (must precede unicode-math) ---
\usepackage{amsmath}
\usepackage{mathtools}

% --- Fonts (lualatex) ---
\usepackage{fontspec}
\usepackage{unicode-math}
\setmathfont{KpMath-Regular.otf}
\setmathfont{KpMath-Bold.otf}[version=bold]

\usepackage{microtype}
\usepackage[dvipsnames]{xcolor}
\usepackage[margin=25mm, top=30mm, bottom=30mm]{geometry}

\usepackage{fancyhdr}
\pagestyle{fancy}
\fancyhf{}
\fancyhead[L]{\small\itshape Polynomial Voxels on an Octree}
\fancyhead[R]{\small\itshape Nestor (2026), draft}
\fancyfoot[C]{\thepage}
\renewcommand{\headrulewidth}{0.4pt}
\renewcommand{\footrulewidth}{0pt}

\setcounter{secnumdepth}{3}
\usepackage{booktabs}
\usepackage{array}
\usepackage{graphicx}
\usepackage{placeins}
\usepackage[numbers,sort&compress]{natbib}
\usepackage[font=small, labelfont=bf, labelsep=period]{caption}
\usepackage{setspace}
\setstretch{1.15}
\setlength{\parskip}{0.4em}
\usepackage[colorlinks=true, linkcolor=blue!60!black,
            citecolor=blue!60!black, urlcolor=blue!60!black]{hyperref}

% --- Macros ---
\newcommand{\bk}{\mathbf{k}}
\newcommand{\bx}{\mathbf{x}}
\newcommand{\bK}{\mathbf{K}}
\newcommand{\bD}{\boldsymbol{\Delta}}
\newcommand{\cE}{\mathcal{E}}
% the differential of every integral and derivative is upright; an italic d is the side of a cell
\newcommand{\dd}{\mathrm{d}}

\title{\bfseries Polynomial Voxels on an Octree\\[4pt]
\large Placing the Cells where the Medium Needs them,\\
with the Error Known before the Solve}
\author{T. M. Nestor\\[2pt]\small Independent researcher, Australia}
\date{Draft, 2 October 2026}

\begin{document}
\maketitle

\begin{abstract}
\noindent
A volume-integral model of seismic waves in a realistically complicated medium cannot afford a uniform
grid: real media vary on localised scales, and a uniform grid spends most of its cells where nothing
varies. We ask whether cubic voxels can be placed where the medium needs them, on an octree, with the
error of the scattering calculation known before it is run. Each leaf of the tree carries the medium's
contrast to polynomial degree $r$ and the wavefield to degree $p$, and couples to every other leaf
through blocks that follow exactly from those between equal cells. We find that the error of the
scattered far field is the sum of two terms, both available before the solve. The first is the error
of the term of first order in the contrast, which needs no solve: each leaf holds the projection of the
incident wave, and for a leaf of uniform contrast its error is a closed-form factor in spherical Bessel
functions, with leading term $-(\bk_{\mathrm{in}}\cdot\bk_{\mathrm{out}})\,d^2/12$ for a constant field
in a cell of side $d$. The second is the nonlinear fraction of the response times the relative
projection error of the medium. On a sphere with a thin graded shell, against its exact solution, the
second term has one constant ($0.015$) for constant and for linear cells, on uniform grids and on
trees, where the raw ratio of error to projection error varies by a factor of two. Refining by both
terms, a tree of linear cells reaches the accuracy of a uniform grid with a quarter of the cells on a
body whose variation fills $3\%$ of its volume. Which basis suits the medium depends on the medium.
Where it varies smoothly the two degrees must be matched, and a constant contrast with a richer field
is worse than a constant field. Where it is piecewise constant on the leaves, a constant contrast is
exact, and a linear field on the tree has $65$ times less error than the uniform grid of constant cells
with the same number of unknowns. This draft records what is established and what is not.
\end{abstract}

\section{Introduction}\label{sec:intro}

\paragraph{The problem this paper serves.} The aim of this series is an efficient model of the seismic
wavefield in a medium of realistic complication. Its method is the volume integral equation of elastic
scattering \citep{Gubernatis1977}, discretised into cubic voxels that interact through the Green's tensor
of a background medium, the multiple-scattering description of \citet{Foldy1945} and \citet{Lax1951}.
The first paper of the series \citep{NestorContinuum2026} determined what a voxel must carry. A voxel
represents two functions, the medium within it and the wavefield within it, each in a basis of its own.
With the wavefield held to Legendre degree $p$ and the contrast of the medium to degree $r$, the scheme
converges at order $\min(2p+2,2r+2)$ in the cell size, and the error that remains belongs to the single
voxel: at second order in the contrast it is the \emph{relative projection error} of the medium, the
fraction of the squared norm of its profile that polynomials of degree $\min(p,r)$ on the cells cannot
represent. The companion paper \citep{NestorGraded2026} builds those voxels in three dimensions.

\paragraph{The question.} Those results are for uniform grids. A real medium is not uniform in its
difficulty: its variation is concentrated at interfaces, inclusions and thin zones, and a uniform grid
fine enough for those is far finer than the rest of the medium needs. The natural cells are then the
leaves of an octree, cubes whose sizes differ by powers of two. This paper asks two things. Can the
cells be placed where the medium needs them? And can the error of the calculation be known before it is
run, from the medium, the tree and the wave alone? The second matters as much as the first. A tree is
built by a rule, the rule needs an indicator of where the error is, and an indicator that needs the
solution is too late.

\paragraph{The answer in brief.} The error of the scattered field separates by its order in the
contrast into two terms, and both are available before the solve (\S\ref{sec:twoterm}). The term of
first order is computed in one pass over the leaves, and for a leaf of uniform contrast it is a closed
form. The rest is the nonlinear fraction of the response times the relative projection error of the
medium, with one constant for every grid and tree tried. A refinement rule built on the two terms
produces trees that a rule built on the medium alone cannot (\S\ref{sec:refine}). And the choice of
basis for the medium follows from the medium (\S\ref{sec:which}): a constant in each leaf is the right
basis where the medium is piecewise constant on the leaves, and the wrong one where it varies smoothly.

\section{Setting}\label{sec:setting}

\paragraph{Leaves.} The body is covered by cubes, the leaves of an octree: leaf $\ell$ has centre
$\bx_\ell$ and half-width $h_\ell$ (side $d_\ell=2h_\ell$), and the half-widths are in ratios that are
powers of two. With $\boldsymbol\xi=(\bx-\bx_\ell)/h_\ell$ the local coordinate, the leaf carries the
nine-component state (displacement and strain) in the Legendre products
$Q_a(\boldsymbol\xi)=P_{a_1}(\xi_1)P_{a_2}(\xi_2)P_{a_3}(\xi_3)$ of total degree $|a|=a_1+a_2+a_3\le p_\ell$,
which is $1$, $4$ or $10$ functions for $p_\ell=0$, $1$ or $2$, and the contrast in the same products to
degree $r_\ell$. The two degrees may differ from leaf to leaf. A leaf with $p=r=0$ is a Haar cell: a
constant medium and a constant field.

\paragraph{The system.} Tested with the same functions, the volume integral equation becomes
\begin{equation}\label{eq:system}
  M_\ell\,\psi_\ell-\sum_{m}\bK(\ell,m)\,E_m\,\psi_m=\langle Q,\psi^0\rangle_\ell ,
\end{equation}
with $\psi_\ell$ the coefficients of the state in leaf $\ell$, $M_\ell$ the Gram matrix of its
functions, $E_m$ the projected contrast of leaf $m$ acting on its field, $\psi^0$ the incident wave and
$\bK(\ell,m)$ the block that couples the two leaves: the Green's tensor integrated against the functions
of both \citep{NestorGraded2026}.

\paragraph{Blocks between leaves of different sizes.} A polynomial on a cube is a polynomial of the
same degree on each of the cube's descendants, so the block between a large leaf and a small one is an
exact, finite sum of blocks between equal cells: the large leaf is replaced by its descendants at the
size of the small one, with its polynomials re-expanded. No integral is needed beyond those between
equal cells. Against six-dimensional quadrature these blocks agree to $10^{-9}$ for leaves one and two
levels apart, and on a tree whose leaves are all equal the solver reproduces the uniform solver to
$10^{-11}$.

\paragraph{Symmetry.} Two pairs of leaves whose relative positions are images of one another under a
reflection or a permutation of the axes have blocks related by the representation of that operation on
the polynomials and on the nine components, whatever the two sizes, because the operation acts on both
cubes alike. One block is therefore computed for each orbit of the cube's group of $48$ operations.
On a tree of $344$ leaves of three sizes this reduced the assembly from $125$\,s to $65$\,s with an
unchanged solution, and made trees with four sizes practicable (\S\ref{sec:status}).

\section{The error in two terms}\label{sec:twoterm}

Write the scattered far field as a series in the contrast: $u=T_1+T_2+\dots$, with $T_1$ the term of
first order (the Born term) and the rest the nonlinear part of the response. The error of the scheme
separates the same way.

\subsection{The first-order error needs no solve}\label{sec:born}

At first order in the contrast Eq.~\eqref{eq:system} is $M_\ell\psi_\ell=\langle Q,\psi^0\rangle_\ell$:
each leaf holds the $L^2$ projection of the incident wave onto its field functions. The scheme's Born
term is therefore known from one pass over the leaves, with no system to solve, and its difference from
the Born term of a finer representation (each leaf's eight children with quadratic cells) is the
first-order error of the tree, leaf by leaf. On a tree of two leaf sizes this projected Born term equals
the weak-contrast limit of the full solver to $10^{-5}$ of its size.

\subsection{The wave factor of a leaf}\label{sec:factor}

For a leaf whose contrast is uniform the first-order error is a closed form. The leaf represents the
incident wave $\mathrm{e}^{\mathrm{i}\bk_{\mathrm{in}}\cdot\bx}$ by its projection, and it radiates
towards the observer through the moments of $\mathrm{e}^{-\mathrm{i}\bk_{\mathrm{out}}\cdot\bx}$, with
$\bk_{\mathrm{out}}$ the P or the S wavevector in the direction of observation. Since
$\int_{-1}^{1}P_n(t)\,\mathrm{e}^{\mathrm{i}qt}\,\dd t=2\,\mathrm{i}^n j_n(q)$, with $j_n$ the spherical
Bessel function, the leaf's Born far field is its exact Born far field times
\begin{equation}\label{eq:factor}
  1+E_p=\frac{\displaystyle\sum_{|a|\le p}\ \prod_{i=1}^{3}(2a_i+1)\,
        j_{a_i}\bigl(k_{\mathrm{in},i}h\bigr)\,j_{a_i}\bigl(k_{\mathrm{out},i}h\bigr)}
       {\displaystyle\prod_{i=1}^{3}j_0\bigl((k_{\mathrm{in},i}-k_{\mathrm{out},i})h\bigr)} .
\end{equation}
The sum over all degrees is the denominator (the addition theorem), so $E_p\to0$ as $p$ grows. Its
leading term is
\begin{equation}\label{eq:lead}
  E_p=-h^{2p+2}\sum_{|a|=p+1}\ \prod_{i=1}^{3}
      \frac{(2a_i+1)\,\bigl(k_{\mathrm{in},i}k_{\mathrm{out},i}\bigr)^{a_i}}{\bigl[(2a_i+1)!!\bigr]^2}
      +O\bigl(h^{2p+4}\bigr),
  \qquad E_0=-\frac{(\bk_{\mathrm{in}}\cdot\bk_{\mathrm{out}})\,d^2}{12}+\dots
\end{equation}
For a constant field the error of a leaf therefore depends on the angle of scattering, vanishes at
right angles to the incident wave, and differs between the P and the S far field. Along one axis, with
$\bk_{\mathrm{out}}=\pm\bk_{\mathrm{in}}$ of magnitude $k$, Eq.~\eqref{eq:lead} is
$-(\pm1)^{p+1}c_p(kd)^{2p+2}$ with $c_p=(2p+3)/\bigl(4^{p+1}[(2p+3)!!]^2\bigr)$: the constants
$c_0=1/12$, $c_1=1/720$ and $c_2=1/100800$ of the layer \citep{NestorContinuum2026}. The layer's
closed-form error is the one-dimensional case of the leaf's.

\paragraph{Checks.} Equation~\eqref{eq:factor} is derived symbolically, built a second way without
Bessel functions (project, multiply, integrate over the cube) and found identical for $p=0$, $1$, $2$;
Eq.~\eqref{eq:lead} and its one-dimensional reduction are verified symbolically; and the implementation
used below agrees with $30$-digit values of the factor to $10^{-13}$ in $36$ cases. Against the solver,
on a homogeneous cube (which every grid of cubes represents exactly) at $k_SA=1$ for half-side $A$: the
Born far field divided by $1+E_p$ agrees across $p=0,1,2$ and two grids to $9\times10^{-7}$, where the
uncorrected fields differ by up to $14\%$; and the Born far field of a tree of $15$ leaves of mixed
sizes and degrees equals the sum over its leaves of exact contribution times factor to $2\times10^{-9}$.

\subsection{The rest is the projection error of the medium}\label{sec:rest}

Let $\cE_r$ be the relative projection error of the medium on the tree,
\begin{equation}\label{eq:proj}
  \cE_r=\frac{\sum_\ell\int_\ell\bigl(f-\Pi_rf\bigr)^2\,\dd V}{\int f^2\,\dd V},
\end{equation}
with $f(\bx)$ the profile of the contrast and $\Pi_r$ the $L^2$ projection onto the polynomials of
degree at most $r$ on each leaf. Subtract from the error of the far field the first-order error of
\S\ref{sec:born}. What is left, divided by $\cE_r$, is one constant (Table~\ref{tab:twoterm}). The test
body is a sphere of radius $a$ with a uniform core of radius $0.75a$ and a shell across which the
contrast falls smoothly to zero, at $k_Sa=0.5$, whose exact solution is known; the error is the largest
difference of the far field from the exact one over nine scattering angles, relative to its peak.
\begin{table}[htbp]
\centering
\caption{The error of the far field in two terms, on the sphere with a graded shell. The rest is the
error less the first-order error, which is computed without a solve. $\cE$ is the relative projection
error of Eq.~\eqref{eq:proj} at the degree of the cells. The nonlinear fraction of the exact response
is $0.0151$.}\label{tab:twoterm}
\small
\begin{tabular}{@{}llrccc@{}}
\toprule
cells & grid & leaves & error & error / $\cE$ & rest / $\cE$ \\
\midrule
constant ($p=r=0$) & uniform & $64$   & $8.17\times10^{-3}$ & $0.0309$ & $0.0149$ \\
 & uniform & $184$  & $3.18\times10^{-3}$ & $0.0235$ & $0.0151$ \\
 & uniform & $408$  & $1.61\times10^{-3}$ & $0.0196$ & $0.0154$ \\
 & uniform & $1256$ & $6.77\times10^{-4}$ & $0.0166$ & $0.0158$ \\
 & tree    & $232$  & $3.29\times10^{-3}$ & $0.0321$ & $0.0152$ \\
 & tree    & $856$  & $9.28\times10^{-4}$ & $0.0210$ & $0.0152$ \\
 & tree    & $1360$ & $6.03\times10^{-4}$ & $0.0210$ & $0.0151$ \\
 & tree    & $1696$ & $5.75\times10^{-4}$ & $0.0222$ & $0.0151$ \\
\midrule
linear ($p=r=1$) & uniform & $64$  & $8.99\times10^{-4}$ & $0.0147$ & $0.0148$ \\
 & uniform & $184$ & $4.62\times10^{-4}$ & $0.0146$ & $0.0146$ \\
 & uniform & $408$ & $2.20\times10^{-4}$ & $0.0146$ & $0.0146$ \\
 & tree    & $184$ & $5.16\times10^{-4}$ & $0.0142$ & $0.0148$ \\
 & tree    & $352$ & $2.20\times10^{-4}$ & $0.0146$ & $0.0145$ \\
\bottomrule
\end{tabular}
\end{table}
For constant cells the ratio of error to projection error varies by a factor of $1.93$ across the
grids; the rest over the projection error varies by $1.06$, and it is $0.0151$ to $0.0152$ on all four
trees. For linear cells the first-order error is small and both ratios are near $0.0146$. The constant
is the nonlinear fraction of the exact response, $\lvert u-T_1\rvert/\lvert u\rvert=0.0151$. This is
the law of the layer \citep{NestorContinuum2026} in three dimensions: the error of the nonlinear part
is that part times the relative projection error of the medium. It holds for constant cells as it does
for linear ones, once the first-order error is accounted for; without that, constant cells appear not
to follow it.

\section{Refinement by the two terms}\label{sec:refine}

\paragraph{The rule.} Each leaf has two indicators, both computed before any solve. The \emph{medium}
indicator is the nonlinear fraction $\nu$ times the leaf's share of the projection error,
$\nu\int_\ell(f-\Pi_rf)^2\,\dd V/\int f^2\,\dd V$. The \emph{wave} indicator is the leaf's own
first-order error (\S\ref{sec:born}) relative to the peak of the Born far field. A leaf is split into
its eight children while the sum of the two exceeds a tolerance. The nonlinear fraction is estimated
from one solve on a coarse tree that resolves the medium, as $\lvert u-T_1\rvert/\lvert u\rvert$ of that
tree. The predicted error of a tree is the first-order error of the whole tree, summed over the leaves
as complex amplitudes, plus $\nu\,\cE_r$. It uses no reference solution.

\paragraph{A localised feature.} The test body is a sphere of radius $a$ whose contrast is a weak
halo, falling smoothly from $5\%$ of the central value to zero across the whole radius, plus a
compact feature in which it falls from the central value to the halo between $0.1a$ and $0.3a$: under
$3\%$ of the volume. Its exact solution is known. At $k_Sa=1$, with leaves no smaller than $a/32$, the
errors are those of Table~\ref{tab:refine}. The rule that uses the medium alone refines a leaf while
its share of the projection error exceeds a tolerance.
\begin{table}[htbp]
\centering
\caption{The localised feature: leaves, the error predicted without a reference, and the error against
the exact solution. The nonlinear fraction estimated from a coarse tree is $0.0051$ (constant cells)
and $0.0059$ (linear cells); the exact value is $0.0059$.}\label{tab:refine}
\small
\begin{tabular}{@{}llrcc@{}}
\toprule
cells & grid & leaves & predicted & error \\
\midrule
constant & uniform & $408$ & $1.11\times10^{-2}$ & $7.15\times10^{-3}$ \\
 & uniform & $1256$ & $4.36\times10^{-3}$ & $2.92\times10^{-3}$ \\
 & medium alone & $344$ & $1.60\times10^{-2}$ & $1.55\times10^{-2}$ \\
 & medium alone & $512$ & $1.57\times10^{-2}$ & $1.53\times10^{-2}$ \\
 & medium alone & $736$ & $5.92\times10^{-3}$ & $5.49\times10^{-3}$ \\
 & two terms & $260$ & $8.18\times10^{-3}$ & $7.18\times10^{-3}$ \\
 & two terms & $372$ & $6.82\times10^{-3}$ & $5.38\times10^{-3}$ \\
 & two terms & $752$ & $4.48\times10^{-3}$ & $3.78\times10^{-3}$ \\
 & two terms & $1780$ & $2.18\times10^{-3}$ & $1.75\times10^{-3}$ \\
\midrule
linear & uniform & $184$ & $2.83\times10^{-3}$ & $2.73\times10^{-3}$ \\
 & uniform & $408$ & $8.27\times10^{-4}$ & $6.77\times10^{-4}$ \\
 & medium alone & $176$ & $2.43\times10^{-4}$ & $2.47\times10^{-4}$ \\
 & medium alone & $344$ & $1.95\times10^{-4}$ & $1.96\times10^{-4}$ \\
 & two terms & $260$ & $1.72\times10^{-4}$ & $1.68\times10^{-4}$ \\
 & two terms & $336$ & $1.56\times10^{-4}$ & $1.66\times10^{-4}$ \\
 & two terms & $452$ & $1.08\times10^{-4}$ & $1.14\times10^{-4}$ \\
\bottomrule
\end{tabular}
\end{table}

\paragraph{Constant cells.} The rule that uses the medium alone fails: with $512$ leaves its error is
twice that of the uniform grid of $408$ cells. It refines the feature and leaves the halo coarse, and
the halo is where the first-order error is, in leaves whose medium is too smooth for the rule to see.
The rule with two terms sees it. Its tree of $260$ leaves has the error of the uniform grid of $408$,
and against the trend of the uniform grids its trees have the same error with about $35\%$ fewer cells
at the coarse end and an error $15$ to $20\%$ lower at the fine end. The saving is modest, and the
reason is the cell: a constant field has a second-order error in every leaf, so the halo needs cells
however smooth its medium is.

\paragraph{Linear cells.} Here the tree pays. The tree of $176$ leaves has eleven times less error than
the uniform grid of $184$ cells. By the trend of the two uniform grids its error would need about $730$
uniform cells, four times as many, and that of the tree of $452$ leaves about $1130$; these uniform
equivalents are extrapolated, not computed. The fourth-order error of a linear field lets the halo stay
coarse, and the tree then follows the medium. The two-term rule improves on the medium alone by about
$15\%$ at about $340$ leaves.

\paragraph{The prediction.} The predicted error is above the true error on all eleven grids of constant
cells computed for this body (nine of them in Table~\ref{tab:refine}), by $3\%$ to $55\%$. For linear
cells the ratio of predicted to true is $0.94$ to $1.22$ on nine grids: close, but not a bound.

\section{Which basis for the medium}\label{sec:which}

\paragraph{A smooth medium: the degrees must be matched.} On the body of \S\ref{sec:refine}, hold the
contrast constant in every leaf and raise the field to linear. The error rises, at four times the
unknowns: from $7.15\times10^{-3}$ to $8.31\times10^{-3}$ on the uniform grid of $408$ cells, and from
$7.18\times10^{-3}$ to $8.32\times10^{-3}$ on the tree of $260$ leaves. The prediction, made without a
solve, says so beforehand ($1.11\times10^{-2}$ against $1.34\times10^{-2}$ on the uniform grid). On a
smoothly varying medium the first-order error of a leaf with a constant contrast is not the wave's: it
is the error of holding the contrast constant, which a richer field cannot repair, and with a constant
field the wave's error partly cancels it. This is the rule of \citet{NestorContinuum2026} that the two
bases must be matched.

\paragraph{A medium that is piecewise constant on the leaves.} The pairing of a constant medium with a
polynomial field belongs to the medium that is constant in each leaf. The test body is a cube of
half-side $B$ with a uniform contrast a quarter of the central value, holding at its centre a cube of
half-side $B/4$ with the full contrast, at $k_SB=2$. A tree of $120$ leaves represents it exactly with
one constant per leaf ($\cE_0=0$): $56$ leaves of half-width $B/4$ and $64$ of half-width $B/8$. The
uniform grid that does the same has $512$ cells. No exact solution is known for this body; the
reference is the tree with quadratic fields, whose own first-order error is predicted to be
$2\times10^{-7}$ (Table~\ref{tab:blocky}).
\begin{table}[htbp]
\centering
\caption{A cube within a cube, the contrast constant in every leaf. The first-order error is predicted
without a solve; the error is against the tree with quadratic fields. In the last two rows the faces of
the inner cube fall inside the cells of the grid.}\label{tab:blocky}
\small
\begin{tabular}{@{}lrrcc@{}}
\toprule
grid and field & cells & unknowns & predicted & error \\
\midrule
tree, constant field & $120$ & $1080$ & $1.54\times10^{-2}$ & $1.53\times10^{-2}$ \\
tree, linear in the $56$ large leaves & $120$ & $2592$ & $3.24\times10^{-3}$ & $3.21\times10^{-3}$ \\
tree, linear field & $120$ & $4320$ & $8.07\times10^{-5}$ & $8.42\times10^{-5}$ \\
uniform, constant field & $512$ & $4608$ & $5.54\times10^{-3}$ & $5.51\times10^{-3}$ \\
uniform, linear field & $512$ & $18432$ & $8.17\times10^{-6}$ & $9.25\times10^{-6}$ \\
\midrule
uniform, constant field & $64$ & $576$ & $1.81\times10^{-2}$ & $1.79\times10^{-2}$ \\
uniform, linear field & $64$ & $2304$ & $3.01\times10^{-2}$ & $2.92\times10^{-2}$ \\
\bottomrule
\end{tabular}
\end{table}
With the medium exact there is no second term, and the first-order prediction is within $13\%$ of the
measured error in all seven cases and within $4\%$ in five. The tree with a linear field has an error
of $8.4\times10^{-5}$ with $4320$ unknowns; the uniform grid of constant cells has $5.5\times10^{-3}$
with $4608$: $65$ times more. Raising the field to linear in the large leaves alone lowers the tree's
error from $1.5\times10^{-2}$ to $3.2\times10^{-3}$: a large leaf needs a richer basis for the
wavefield, not for the medium. And where the faces of the inner cube fall inside cells, the linear
field is worse than the constant one, as on the smooth body.

\paragraph{The choice.} The two cases give the rule. Where the medium is piecewise constant and the
tree can follow its interfaces, hold it constant in each leaf and spend the degree on the field.
Where it varies smoothly, give the medium and the field the same degree and let the tree follow the
medium. In both the error is predicted before the solve.

\section{What is established, and what is not}\label{sec:status}

\begin{itemize}
\item \emph{Established twice.} The wave factor, Eqs.~\eqref{eq:factor} and \eqref{eq:lead}: a symbolic
derivation and an independent numerical implementation, and both against the solver.
\item \emph{Established against exact solutions, by one implementation.} The two-term separation
(Table~\ref{tab:twoterm}) and the refinement results (Table~\ref{tab:refine}): one frequency each, two
bodies, both spheres with radial profiles.
\item \emph{Self-consistent only.} The cube within a cube (Table~\ref{tab:blocky}) has no exact
solution; its reference is the same scheme with quadratic fields. The field near the edges and corners
of the cubes is not smooth, and an independent solution is needed before the figure of $65$ is relied
on.
\item \emph{Not yet derived.} The nonlinear fraction $\nu$ is estimated from a coarse solve; it is
within $20\%$ of the exact value in the cases tried. The constant of the second term is observed to
equal the nonlinear fraction of the response; the layer's derivation \citep{NestorContinuum2026} has
not been carried to three dimensions.
\item \emph{Not a bound.} The predicted error exceeded the true error for constant cells in every case,
but fell $6\%$ below it for linear cells.
\item \emph{Size.} The solver is dense and stops near $20\,000$ unknowns: about $2200$ constant leaves
or $550$ linear ones. The saving in cells is therefore shown on small bodies. An equal-cell block costs
about $50$\,ms, and a tree of $344$ leaves of four sizes is assembled and solved in about three
minutes.
\item \emph{Earlier work.} The relation of these results to adaptive methods for integral equations and
to the fast multipole method on an octree is not yet written; the literature has not been surveyed.
\end{itemize}

\section{Conclusions}\label{sec:summary}

On an octree of polynomial voxels the error of the scattered field is two terms, and both are known
before the solve: the first-order error, from one pass over the leaves and in closed form for a leaf of
uniform contrast, and the nonlinear fraction of the response times the relative projection error of the
medium. The closed form reduces in one dimension to the constants of the layer. A refinement rule built
on both terms places cells where a rule built on the medium alone cannot, and with linear cells a tree
reaches the accuracy of a uniform grid with a quarter of the cells on a body whose variation is
localised. A constant contrast in each leaf is the right basis for a medium that is piecewise constant
on the leaves, where a linear field on the tree outperforms the uniform grid of constant cells
sixty-five-fold at equal unknowns in a self-consistent test, and the wrong one for a smooth medium,
where the two degrees must be matched. What remains before the method models a realistic medium is a
solver whose cost does not grow as the square of the number of leaves.

\vfill
\noindent\rule{\linewidth}{0.4pt}\\
{\footnotesize The octree, its blocks and its solver are \path{cubic_scattering/graded_voxel/octree.py},
with tests in \path{cubic_scattering/tests/test_graded_voxel_octree.py}. The wave factor is derived in
\path{Mathematica/AdaptiveOctree_WaveFactor.wl} and measured by
\path{scripts/measure_octree_wave_term.py}; Table~\ref{tab:twoterm} is
\path{scripts/measure_octree_two_term.py}; Table~\ref{tab:refine} is
\path{scripts/pilot_octree_two_term_refinement.py}; the constant medium with a richer field is
\path{scripts/pilot_octree_hp_refinement.py}; Table~\ref{tab:blocky} is
\path{scripts/pilot_octree_blocky_body.py}. Their results are stored in
\path{LatexPDFs/AdaptiveOctree/figures}.}

\bibliographystyle{unsrtnat}
\bibliography{references}

\end{document}
